\begin{remark*}[Geometric intuition]
A right triangle fixes three side lengths relative to an acute angle:
opposite, adjacent, and hypotenuse. Similar right triangles may be larger or
smaller, but their corresponding side ratios are unchanged. This invariance is
what makes a trigonometric ratio a function of the angle rather than a feature
of one drawing.
\end{remark*}

\begin{center}
\begin{tikzpicture}[scale=1.25, line cap=round, line join=round]
  \definecolor{trigblue}{RGB}{30,110,190}
  \coordinate (A) at (0,0);
  \coordinate (B) at (3.4,0);
  \coordinate (C) at (3.4,2.0);
  \draw[very thick] (A)--(B)--(C)--cycle;
  \draw (3.1,0)--(3.1,0.3)--(3.4,0.3);
  \draw[trigblue] (0.55,0) arc[start angle=0,end angle=30.5,radius=0.55];
  \node[trigblue] at (0.78,0.22) {$\theta$};
  \node[below] at ($(A)!0.5!(B)$) {adjacent};
  \node[right] at ($(B)!0.5!(C)$) {opposite};
  \node[above left] at ($(A)!0.5!(C)$) {hypotenuse};
  \foreach \p/\pos in {A/below left,B/below right,C/above right}{
    \fill (\p) circle (0.025);
    \node[\pos] at (\p) {$\p$};
  }
\end{tikzpicture}
\end{center}


\begin{remark*}[Primitive and derived ratios]
The primitive geometric ratios are sine and cosine. Tangent, cotangent, secant,
and cosecant are convenient quotient or reciprocal ratios built from the same
three side lengths.
\end{remark*}
\begin{dependencies}
\begin{itemize}
  \item \hyperref[def:euclid-triangle-angle-classes]{Right-angled triangle}
  \item \hyperref[def:euclid-right-angle-and-perpendicular]{Right angle and perpendicular}
\end{itemize}
\end{dependencies}
